\section*{Chapitre \ref{ch:g11:infrared} --- La spectroscopie infrarouge}

\begin{solution}{exo:g11:infrared:1}
$\lambda = 1/3300 = \qty{3.03e-4}{cm} = \qty{3.03}{\micro m}$~;
$\lambda = 1/1050 = \qty{9.52e-4}{cm} = \qty{9.52}{\micro m}$. La bande à
\qty{3300}{cm^{-1}}, de \omterm{def:g11:infrared:wavenumber}{nombre d'onde} plus grand, correspond à la
vibration la plus rapide.
\end{solution}

\begin{solution}{exo:g11:infrared:2}
$\qty{10}{\micro m} = \qty{1.0e-3}{cm}$, donc $\sigma = 1/\num{1.0e-3} =
\qty{1000}{cm^{-1}}$.
\end{solution}

\begin{solution}{exo:g11:infrared:3}
Le \ce{C=O} d'une \omterm{def:g11:functional-groups:carbonyl}{cétone}~; un \ce{O-H} d'\omterm{def:g11:functional-groups:alcohol}{alcool} lié par \omterm{def:g11:polarity-and-cohesion:hydrogen-bond}{liaison
hydrogène}~; les liaisons \ce{C-H} d'une chaîne carbonée.
\end{solution}

\begin{solution}{exo:g11:infrared:4}
La bande à \qty{1715}{cm^{-1}}~: la liaison \ce{C=O}.
\end{solution}

\begin{solution}{exo:g11:infrared:5}
L'instrument mesure la lumière qui traverse l'échantillon~; là où une
liaison absorbe, moins de lumière passe et la courbe plonge. $T =
\qty{100}{\percent}$ signifie qu'aucune lumière de ce \omterm{def:g11:infrared:wavenumber}{nombre d'onde} n'est
absorbée.
\end{solution}

\begin{solution}{exo:g11:infrared:6}
Le premier a une bande \ce{C=O} à \qty{1715}{cm^{-1}} et pas de bande
\ce{O-H}~: c'est une \omterm{def:g11:functional-groups:carbonyl}{cétone}, la butanone (ou butan-2-one),
\ce{CH3-CO-CH2-CH3}. Le second a une bande \ce{O-H} et pas de
\ce{C=O}~: c'est un \omterm{def:g11:functional-groups:alcohol}{alcool}, par exemple le butan-1-ol.
\end{solution}

\begin{solution}{exo:g11:infrared:7}
Les groupes \ce{O-H} de l'acide sont très fortement liés par \omterm{def:g11:polarity-and-cohesion:hydrogen-bond}{liaisons
hydrogène} (les \omterm{def:g7:atoms-and-molecules:molecule}{molécules} s'associent par paires, chaque \ce{O-H} étant
lié au \ce{C=O} de l'autre), et chaque \omterm{def:g7:atoms-and-molecules:molecule}{molécule} est liée un peu
différemment~: la liaison est encore plus affaiblie et la bande s'étale
sur un domaine très large.
\end{solution}

\begin{solution}{exo:g11:infrared:8}
Dans le liquide (A), la bande \ce{O-H} est large, vers
\qty{3350}{cm^{-1}}~; dans le gaz (B), où les \omterm{def:g7:atoms-and-molecules:molecule}{molécules} sont éloignées et
ne forment pas de \omterm{def:g11:polarity-and-cohesion:hydrogen-bond}{liaisons hydrogène}, c'est une bande fine vers
\qty{3666}{cm^{-1}}.
\end{solution}

\begin{solution}{exo:g11:infrared:9}
Difficilement~: vers \qty{1730}{cm^{-1}} pour l'\omterm{def:g11:functional-groups:carbonyl}{aldéhyde} et
\qty{1715}{cm^{-1}} pour la \omterm{def:g11:functional-groups:carbonyl}{cétone}, des valeurs proches. Comparer tout le
spectre, \omterm{def:g11:infrared:band}{empreinte digitale} comprise, à des spectres de référence
permettrait de trancher, ou une autre technique.
\end{solution}

\begin{solution}{exo:g11:infrared:10}
Un \omterm{def:g11:functional-groups:carboxylic-acid}{ester} (\ce{C=O} vers \qty{1735}{cm^{-1}}, \ce{C-O} intense, pas de
\ce{O-H}), par exemple l'éthanoate d'éthyle. L'acide éthanoïque a pour
formule \ce{C2H4O2}, et non \ce{C4H8O2}. L'acide butanoïque a la bonne
formule mais présenterait la bande \ce{O-H} très large des acides~: il
est exclu.
\end{solution}

\begin{solution}{exo:g11:infrared:11}
Le propan-1-ol est un \omterm{def:g11:functional-groups:alcohol}{alcool}~: comme A. L'acide propanoïque est un \omterm{def:g11:functional-groups:carboxylic-acid}{acide
carboxylique}~: comme D.
\end{solution}

\begin{solution}{exo:g11:infrared:12}
La large bande \ce{O-H} vers \qty{3350}{cm^{-1}} diminue tandis qu'une
bande \ce{C=O} intense vers \qty{1715}{cm^{-1}} grandit. La réaction est
terminée quand la bande \ce{O-H} a disparu et que la bande \ce{C=O} ne
grandit plus.
\end{solution}

\begin{solution}{exo:g11:infrared:13}
L'acide butanoïque présente une bande \ce{O-H} très large de 2500 à
\qty{3300}{cm^{-1}}, l'\omterm{def:g11:functional-groups:carboxylic-acid}{ester} aucune~; le \ce{C=O} de l'acide se situe
vers \qty{1710}{cm^{-1}}, plus large, celui de l'\omterm{def:g11:functional-groups:carboxylic-acid}{ester} vers
\qty{1735}{cm^{-1}}, fin, avec une bande \ce{C-O} intense vers
\qty{1240}{cm^{-1}}. L'acide possède un groupe \ce{O-H}, lié par \omterm{def:g11:polarity-and-cohesion:hydrogen-bond}{liaison
hydrogène}, qui affaiblit aussi légèrement son \ce{C=O}.
\end{solution}

\begin{solution}{exo:g11:infrared:14}
De l'éthanol restant (l'\omterm{def:g11:functional-groups:alcohol}{alcool} à partir duquel l'\omterm{def:g11:functional-groups:carboxylic-acid}{ester} est fabriqué) ou
de l'eau (venue de l'\omterm{def:g4:air-a-mixture-of-gases:air}{air} ou du lavage)~: tous deux ont des groupes
\ce{O-H} liés par \omterm{def:g11:polarity-and-cohesion:hydrogen-bond}{liaison hydrogène}. On peut comparer l'échantillon à un
spectre de référence, le sécher, le distiller, ou mesurer sa température
d'ébullition.
\end{solution}

\begin{solution}{exo:g11:infrared:15}
Une \omterm{def:g10:lewis-and-shape:multiple-bond}{liaison double} est un ressort plus raide qu'une liaison simple entre
les mêmes \omterm{def:g7:atoms-and-molecules:atom}{atomes}~: elle vibre plus vite, à un \omterm{def:g11:infrared:wavenumber}{nombre d'onde} plus élevé
(1715 contre \qty{1050}{cm^{-1}}). Les liaisons \ce{O-H}, \ce{N-H} et
\ce{C-H} font toutes intervenir un \omterm{def:g7:atoms-and-molecules:atom}{atome} d'hydrogène, le plus léger de
tous~: des boules légères sur un ressort vibrent vite, d'où des \omterm{def:g11:infrared:wavenumber}{nombres
d'onde} vers 2900--\qty{3600}{cm^{-1}}.
\end{solution}

\begin{solution}{pb:g11:infrared:1}
\textbf{1.} \ce{CH3-CH2-OH}~; \ce{CH3-CO-CH3}~; \ce{CH3-COOH}.

\textbf{2.} \ce{C2H6O}~; \ce{C3H6O}~; \ce{C2H4O2}.

\textbf{3.} Hydroxyle~; carbonyle (\omterm{def:g11:functional-groups:carbonyl}{cétone})~; carboxyle.

\textbf{4.} L'éthanol~: \ce{O-H} et \ce{C-H}. La propanone~: \ce{C-H} et
\ce{C=O}. L'acide éthanoïque~: \ce{O-H}, \ce{C-H} et \ce{C=O}.

\textbf{5.} Pas de bande \ce{C=O}, une large bande \ce{O-H}~: l'éthanol.

\textbf{6.} Une bande \ce{O-H} très large de 2500 à \qty{3300}{cm^{-1}}
et un \ce{C=O} intense vers \qty{1710}{cm^{-1}}~: l'acide éthanoïque.

\textbf{7.} La propanone~: la bande \ce{C=O} très intense à
\qty{1715}{cm^{-1}}, sans \ce{O-H}.

\textbf{8.} Probablement (vinaigre, \omterm{def:g6:solutions-and-solubility:solute}{solvant}, \omterm{def:g11:functional-groups:alcohol}{alcool}), mais renifler des
liquides inconnus est dangereux~; le spectre est sûr et certain.

\textbf{9.} Les groupes \ce{O-H} de l'acide sont plus fortement liés par
\omterm{def:g11:polarity-and-cohesion:hydrogen-bond}{liaisons hydrogène} (\omterm{def:g7:atoms-and-molecules:molecule}{molécules} associées par paires)~: la bande est plus
basse et beaucoup plus large.

\textbf{10.} \ce{CH3-O-CH3}, \ce{C2H6O}~; il n'a pas de \ce{O-H}, donc
pas de bande \ce{O-H}.

\textbf{11.} La bande \ce{C=O}~: vers \qty{1730}{cm^{-1}} pour le
propanal, vers \qty{1715}{cm^{-1}} pour la propanone.

\textbf{12.} La bande \ce{O-H}~: un \omterm{def:g11:functional-groups:carboxylic-acid}{ester} n'a pas de \ce{O-H}. Son
\ce{C=O} se situerait vers \qty{1735}{cm^{-1}}.

\textbf{13.} Des \omterm{def:g11:organic-skeletons:isomer}{isomères} ont la même formule~; leurs \omterm{def:g11:functional-groups:functional-group}{groupes
caractéristiques} diffèrent, et le spectre montre les groupes.

\textbf{14.} \qty{1715}{cm^{-1}}.

\textbf{15.} $1715 \times 100 = \qty{1.715e5}{m^{-1}}$.

\textbf{16.} $\lambda = 1/\num{1.715e5} = \qty{5.83e-6}{m} =
\qty{5.83}{\micro m}$.

\textbf{17.} Non~: la lumière visible va d'environ 0.4 à
\qty{0.75}{\micro m}. C'est un rayonnement infrarouge.

\textbf{18.} $\num{1.05e5}\ \si{m^{-1}}$, donc $\lambda = \qty{9.52e-6}{m}
= \qty{9.52}{\micro m}$.

\textbf{19.} La liaison \ce{C=O} (\omterm{def:g11:infrared:wavenumber}{nombre d'onde} plus élevé, longueur
d'onde plus courte).

\textbf{20.} Environ \qty{5.8}{\micro m}.
\end{solution}
